\section{Introduction}
\label{sec:intro}

When humans reason about complex problems, they rarely enumerate numbered steps on a mental whiteboard. Instead, they often construct narratives: a student imagines a shopkeeper making change to solve an arithmetic problem; an engineer mentally simulates a bridge under load to predict failure modes. This narrative mode of reasoning---building stories with characters, settings, and causal progressions---is well-documented in cognitive science \citep{bruner1991narrative} and stands in contrast to the formulaic step-by-step decomposition that has become the default in \chainofthought prompting for large language models \citep{wei2022chain}.

\para{The appeal of narrative reasoning.} \chainofthought prompting dramatically improves LLM performance on tasks requiring multi-step reasoning \citep{wei2022chain, wang2023selfconsistency}. Yet the standard format---``Step 1: \ldots Step 2: \ldots''---is only one of many possible reasoning structures. Recent work on Narrative-of-Thought \citep{zhang2024narrative} demonstrated that prompting LLMs to reason through temporally grounded narratives yields 16--71\% F1 improvements over standard prompting on temporal graph generation tasks. This raises a natural question: does narrative-form CoT generalize to other reasoning domains?

\para{Our main question.} {\bf Does prompting LLMs to reason in story form improve, match, or degrade performance compared to standard step-by-step CoT?} If narrative structure activates different model capabilities---perhaps leveraging the vast narrative text in LLM training corpora---it could unlock improvements across diverse tasks. Alternatively, if the benefit of CoT is purely computational scaffolding \citep{schaeffer2023invalid}, then format should not matter.

\para{What we do.} We introduce \storycot, a prompting strategy in which few-shot exemplars wrap identical logical steps inside vivid narratives with characters, settings, and temporal progression. We systematically compare four prompting strategies---\directprompt, \zeroshotcot, \fewshotcot, and \storycot---across four benchmarks spanning math (\gsm), commonsense (\csqa), multi-hop (\strategyqa), and science (\arc) reasoning, all using \gptfour.

\para{What we find.} \storycot performs equivalently to \fewshotcot across all four domains. Accuracy differences range from 0 to 1.3 percentage points, and none reach statistical significance (McNemar's test, all $p > 0.47$). Both structured methods dramatically outperform \directprompt on multi-step tasks (+36 pp on \gsm, +19 pp on \arc). This null result is itself informative: it suggests that for state-of-the-art LLMs, \chainofthought benefits arise from the \emph{presence} of intermediate reasoning tokens, not from their format.

Our contributions are:
\begin{itemize}[leftmargin=*,itemsep=0pt,topsep=0pt]
\item We conduct the first systematic comparison of narrative-form \chainofthought against standard \chainofthought across four diverse reasoning benchmarks, finding format equivalence on all tasks.
\item We provide evidence for the \emph{computational scaffolding hypothesis}: CoT improves LLM reasoning through intermediate computation, independent of whether that computation takes the form of logical steps or narrative prose.
\item We show that \zeroshotcot can \emph{harm} performance on commonsense reasoning (--14 pp on \csqa), highlighting that the structure of CoT prompts matters even when their format does not.
\item We release our prompts, evaluation code, and full results to support future work on alternative reasoning formats.
\end{itemize}
